\documentclass[12pt]{book}

\usepackage{graphicx}
\usepackage{eso-pic}
\usepackage{tcolorbox}
\usepackage{polyglossia}
\setdefaultlanguage[indentfirst=false,forceheadingpunctuation=false]{russian}
\setotherlanguages{english}

\RequirePackage[a4paper, headsep = 0.5 \headsep, left=2.5cm, right=2.1cm, top=2cm, bottom=2.1cm]{geometry}

\setmainfont{URW Bookman}
\newfontfamily\cyrillicfont{URW Bookman}[Script=Cyrillic, Ligatures=TeX]
\setlength{\parindent}{20pt}

\usepackage{enumitem}
\setlist[itemize, 1]{label =\raisebox{-0.3\height}{\scalebox{1.8}{\textbullet}}}

\title{Феноменология чуждого}
\author{Бернхард Вальденфельс}
\date{2006}
\begin{document}

\AddToShipoutPictureBG*{\includegraphics[width=\paperwidth]{title.png}}

\null
    \thispagestyle{empty}%
    \addtocounter{page}{-1}%
    \newpage

\tableofcontents

\chapter{Грани чуждого}

Выделить чуждое как отдельную тему --- значит уже упустить его, начать со знакомого места и планировать к нему же вернуться. Обращение к чуждому наверняка приведет к переменам и даже катарсису, хотя в конечном счете привычность и возобладает и может даже углубиться. Чуждое не безобидно и может отдалить нас от самих себя --- потому ему стараются сопротивляться, избегать его или ассимилировать. Эти старания оставляют субъекта наедине с самим собой и требуют выстраивать и оборонять порядок, исключающий проявление чуждого изнутри. Чуждое может пробудить любопытство и рассказать нам о нас самих, но как только оно подходит к тайне свободы и разума --- срабатывает сигнал ''Хаос!'', и свобода и разум остаются в своих правах. Иначе им пришлось бы отказаться от самих себя. Чуждость неизбежно ведет к враждебности, обе стороны все больше убеждаются в своей правоте. Мы становимся ближе и бдительней друг к другу. Встроенные в эмпирическую сеть защитные механизмы исходят из \textit{своего} и стремятся закрепиться в \textit{общем}. Чуждое предполагается внешним, ожидается, что оно будет идентифицировать себя и выглядеть незваным гостем. Его тогда можно будет подвергнуть оценке и суждению и поместить на повседневный моральный, политический, религиозный, культурный и интеллектуальный карантин.

Если же чуждое воспринимать как неизбежно неопределенное, как всегда сокрытое в \textit{своем}, где оно тревожит, соблазняет или пугает, ускользает от понимания, то опыт чуждого, значит, воздействует на \textit{наш} опыт, становясь отчуждением от опыта. Чуждость самореферентна и заразительна, ее последствия предшествуют всякой ее тематизации. Феноменология чуждого, предвосхищенная Гуссерлем в нескольких его новаторских концепциях (чуждый опыт, чуждое эго, чуждое тело и чуждый мир) обладает задатками науки, могущей говорить о чуждом во всей его полноте, всерьез принимать его радикальность, не довольствуясь сведением к его пониманию и практическому распознаванию. Феноменолог, всерьез принимающий эти требования ксенологии, наверняка будет вдохновляться работами Марселя Мосса, Георга Зиммеля, Вальтера Беньямина, Людвига Витгенштейна или Михаила Бахтина, не говоря уже о таких литераторах как Кафка, Музиль, Целан, Валери и Кальвино. В мире искусства другие порядки, и художественное чутье легко преодолевает границы самоуверенности и принципиальности, если вступает в контакт с чуждым.

Настоящее исследование не пытается охватить проблематику чуждого во всей ее пугающей полноте, а лишь обозначает некоторые основные мотивы, побудившие автора разработать феноменологию чуждого в соответствующей серии книг. Первый том --- \textit{Жало чуждого} вышел в 1990г.\footnote{Формальное изложение феноменологии чуждого можно найти в Главе 1 \textit{Топографии чуждого} (\textit{Topographie des Fremden}, 1997). Аспекты чуждого рассматриваются в главах 5-6 книги \textit{Рваные линии опыта} (\textit{Bruchlinien der Erfahrung}, 2002).} Мотивы здесь развивают тему, а не просто ретроактивно обосновывают обращение к ней. Опыт чуждого --- как и опыт удивления, ужаса, любви, ненависти или тех едва заметных несостыковок, через которые возникнет новое --- нельзя вполне описать или обосновать. Разделение на главы выполнено так, что каждая рассматривает конкретный программный компонент проблематики чуждого. Они выстроены в порядке этапов формирования феноменологии чуждого --- они потому вводные, поясняющие тексты и дополняют уже опубликованные обширные исследования. В последующих абзацах я свяжу эту работу с предыдущими, что избавит меня от необходимости отдельно на них ссылаться и позволит сосредоточиться на ключевых темах: порядок -- пафос -- ответ -- тело -- внимание -- межкультурность.

В главе 1 рассматривается \textit{Порядок в сумерках} (\textit{Ordnung im Zwielicht}, 1987). На границах всякого порядка чуждое предстает экстраординарным, не находящим места в порядке --- и все же оно, будучи исключаемым, не ничто. Поскольку оно исключается не вообще --- а из конкретного порядка, то оно нечто большее, чем просто сливающееся с фоном в неопределенности. И поскольку оно взывает к нашим чувствам, то оно не просто шум, от которого мы хотели бы укрыться. Положение, в котором мы оказываемся, обнаруживая несхватываемое в схватываемом, беспорядочное в упорядоченном, невидимое в видимом, тишину в озвучиваемом, запускает цепь исследований, сводящихся к вопросу: как мы можем взаимодействовать с чуждым без предварительной нейтрализации или отрицания его эффектов, его вызовов и требований? Как феноменолог, я всегда старался избегать того, в чем нетерпеливая аудитория любит обвинять феноменологов чуждого --- возведения его в ранг некоего морального или религиозного авторитета или ''совершенно иного''. Это было бы обожествлением Другого --- в чем часто обвиняют Левинаса --- и морализацией чуждого. Но навязать переживание столь же возможно, как и, согласно Канту, чувство. Не то чтобы восприятие чуждого было лишь субъективным ощущением, как в том убеждает модерн, просто в этом мире вещей гораздо больше, чем может охватить школа мысли с ее нормами, ценностями и интересами.

К переживаниям можно относиться с разной степенью осторожности: позволить им говорить самим за себя, использовать как информационный ресурс или пытаться исключить --- будь то от небрежности, властолюбия или безразличия. В главах 2 и 3 эти возможности рассмотрены с точки зрения ответовой этики --- этики, возникающей в ответе, отклике, реакции, и так остающейся за рамками заповедей и запретов. До-предикатный опыт в гуссерлевском смысле здесь обращается в до-нормативный --- остающийся всегда позади, пока мы намечаем действия и строим планы, требующие моральной или правовой оценки. В главе 2 я буду ссылаться в основном на \textit{Рваные линии опыта} (\textit{Bruchlinien der Erfahrung}, 2002): анализ, проведенный в этом тексте, построен вокруг двойного события, состоящего из пафоса и последующей реакции. В начале есть то, что предстает как события --- происходящим с нами, в чем я или Другой можем участвовать --- а не как волящий и действующий субъект. Ключевое здесь --- временной сдвиг. Не нечто иное предшествует инициативе ''я'' --- эта инициатива предшествует самой себе, не возникая из неоткуда. Эти связи предстанут в новом свете, когда в главе 3, рассматривая \textit{Регистр ответов} (\textit{Antwortregister}, 1994)\footnote{В данном случае я намеренно отклоняюсь от порядка, в котором были изложены эти идеи, чтобы облегчить понимания, поскольку \textit{Регистр ответов} предполагает некоторую часть того, что полностью раскрывается только в \textit{Рваных линиях опыта}. ''На что'' ответа указывает на ''чем'' преодоления.}, я раскрываю основные черты ответового поведения и возникающую в результате зова чуждого логику ответов. Парадокс творческого ответа представляет особый случай: оказывается, мы можем воспринимать ''что'' ответа как нечто изобретенное, но не можем сделать то же с его ''на что''. Зов чуждого может быть столь же слабо связан с этическим фундаментализмом, сколь и с технологическим конструктивизмом.

Глава 4 посвящена телесному опыту, и не в коем случае не просто расширяет сферу чуждого. В узле самореференции и самоустранения, определяющих его способом существования, тело оказывается как бы эмблемой чуждого. Чужеродность начинается не с меня самого --- это мало что значило бы,--- так называемый субъект выпадает из центра, когда предстает как \textit{Воплощенное Я} (\textit{Leibliches Selbst})\footnote{Название лекций о теле, опубликованых в 2000г., в которых эта тема рассматривается во всей ее глубине.}. В главе 5 уточняются и подкрепляются ссылками на литературу основные идеи \textit{Феноменологии внимания} (\textit{Phänomenologie der Aufmerksamkeit}, 2004). Внимание полностью подчинено двойному ритму становления очевидным и становления притягательным, двойное движению пафоса и ответа потому преобретает форму, пронизывающую весь наш опыт и поведение. Что становится очевидным или мыслимым, что доходит до наших глаз или ушей, до наших органов чувств, что до ходит до нас --- не наше. Аллюзия названия \textit{Пороги смысла} (Sinnesschwellen, 1999) указывает, что вездесущее внимание ведет в область эстетики и особого этоса чувств. В последних двух главах говорится также о чуждых аспектах психоанализа и технологии, что занимало значительное место в \textit{Рваных линиях опыта} и \textit{Феноменологии внимания}. Как изначальный опыт пронизан моментами неизведанного и немотивированного, так и феномены в нем с самого начала отмечены подлинной феноменотехнологией и соматотехникой. Технология очень долго --- вопреки всем сетованиям современности об отчуждении в ней --- рассматривалась как расширение сферы собственности, но теперь все больше становится источником особой отчужденности.

Главу 6 можно читать как обзор, собрание различных аспектов чуждого в кратком, ориентированном на отношения между культурами, изложении. Межкультурность не ограничена философией.\footnote{Позвольте здесь упомянуть о выходящем в Вене журнале по межкультурной философии Polylog, в котором к настоящему времени появилось довольно много тематических выпусков.} Глобальное мышление не следует ни ожидать, ни желать --- оно неизбежно сформировало бы лишь новые гегемонистские отношения, разве что другими средствами. Философия, переставшая вращаться вокруг себя самой и пустившая корни в привычной культуре, дала бы облегчение. Попытка пересечь границы, не уничтожив их --- одно из будущих начинаний в области межкультурного взаимодействия, и оно значит нечто большее, чем просто сопоставление или смешение культур. Разнообразие слишком легко становится накоплением --- спокойным перебором любопытных вариаций. Опасность для культур, достигних зрелости, заключается в обращении к обществу охотников и собирателей.

Настоящая работа основана на серии отдельных исследований, появившихся в связи с разными обстоятельствами. Все эти исследования были переработаны и ориентированы на общую тему.\footnote{Эти идеи рассматриваются в \textit{Границы нормализации} (\textit{Grenzen der Normalisierung}, 1998).} Я постарался избежать ретроактивной систематизации и сохранить эвристическую множественность подходов, в то же время создав рельефы, благодаря которым --- как на объемной открытке --- были бы видны все значимые аспекты. Как писатели, мы иногда сожалеем, что письмо ограничено двухмерной плоскостью, но, к счастью, в письменной речи --- как и в чтении --- есть свои способы расставления акцентов, и эти подъемы и спуски не дают определенному направлению письма свестись к написанному.

\chapter{Человек как пограничное существо}

Чуждое --- это предельное как оно есть. Даже возникая в знакомом, оно приходит извне. Не бывает чуждого без чуждого места, и порядок нашей жизни, опыта, языка и действий определяет, сколь значимо для нас будет чуждое. Преобразование порядка значит и преобразование чуждого --- такого же многогранного, как и соответствующий порядок. О чуждом говорят не меньше, чем об эго --- пограничные зоны между и за порядками питают чуждое.

\section{Порядки и их пределы}

Обычно верно полагать, что не только порядок ограничен, но и что эти границы возникают в процессе упорядочивания. Нечто --- то, что оно есть в силу его отделения от прочих вещей: камней, растений, животных, природных или искусственных объектов. Однако, в мире людей понятие границ вызывает особое беспокойство, потому что границы всегда под сомнением. Поведение человеческого существа не определяется ни инстинктами, ни искусственными программами: оно не ограничено, а, скорее, особым образом соотнесено с границами. Это касается и места и времени --- границ, определяющих что такое здесь и сейчас,--- и запретов, сдерживающих желания и поступки, и пределов понимания, сдерживающих мышление. Неудивительно, что вопрос о границах существования остается в числе главных тем на протяжении всей человеческой истории: Иегова отделяет свет от тьмы, границы души неизмеримы, философ намечает границы чистого разума, системный теоретик воплощает возвышенный жест созидания в простом ''Проведи границу!''.

Видно, что установление границ, приводя к разным порядкам и структурам, носит не только прагматический и региональный, но и эпохальный характер. Каждая эпоха --- или каждая культура, общество, среда или форма жизни --- ведет себя в рамках некоторых границ, и отношение к этим границам, всегда сопровождаясь определенной политикой, значительно варьируется. Обращение с границами прямо показывает глубинный дух эпохи и может также послужить комментарием к тому, что так долго продвигало современность вперед или предшествовало ей, или подрывало ее и выходило за ее пределы. Рубеж к 21-ому столетию --- определенно очень специфическая граница, и у нас нет языка, подходящего для описания того, что ждет нас впереди.

\section{Безграничная вселенная}

Для контраста можно начать с безграничной вселенной. В нашей культурной традиции она наиболее емко представлена греческим космосом. Космос долгое время имел парадигмальное значение, олицетворяя классическую форму порядка. Он воплощал не один из множества порядков, а порядок вообще, как таковой, и потому единственной альтернативой ему были неупорядоченное многообразие хаоса. Этот космос придает каждому существу его собственную форму (πέρας), границы которой отделяют существо от себя и среды. Точная форма выражается в концептуальном определении (όρισμός); со времен Платона диалектика связывала каждое бытие как единое с прочим в цепочке отношений. Эта горизонтальная связь дополняется вертикальной иерархией, пропорции которой соответствуют степени отражения рационального целого в конкретном существе. Люди здесь стоят выше животных, греки --- выше варваров, мужчина --- выше женщины, созерцание --- выше действия. Рациональность, раскрывая закон целого, имеет решающее значение для места человека в иерархии существ.

Эта реляционная структура --- все различия в ней относительны --- имеет нижнюю и верхнюю границы. Нижняя формируется \textit{ин-дивидом}\footnote{in-dividual; англ. divide --- разделять.}, άτομον είδος --- он уже не может быть далее разделен без разрушения его идентичности. Верхняя граница --- Вселенная, не упорядоченная в рамках чего-либо другого; мир --- сущность, из которой нечего не исключено, постановит Аристотель в своей \textit{Физике} (2.6.207a8). Проще говоря, космос --- порядок без внешнего, с только лишь внутренними границами. И пересекающий эти границы либо попадает в дурную бесконечность бездонного и бессмысленного апейрона, либо, подобно Икару, поднявшемуся на запредельные высоты, терпит фатальную неудачу.

Но замкнутый всеобъемлющий космос, в котором есть место каждому, предполагает, что место проявления и выражения целого как \textit{целого} находится внутри самого этого целого. Психе --- согласно Аристотелю, ''в некотором смысле --- все'' --- становится ступенью порядка, к которому мимически приспосабливается. Космос потому предстает порядком, раскрывающим и выражающим \textit{себя как себя} через превращение всякого условия в момент себя. Такой порядок без внешнего соответствует \textit{мысли о внутреннем} --- \textit{penser du dedans}, перефразируя название работы Фуко\footnote{La Pensée du dehors, 1966} --- и в целом такая мысль была бы мыслью самой по себе. Однако даже древние греки знали о выпадающих из этой схемы фигурах: Сократе --- άτοπος, живущем в городе подобно прочим гражданам, но действующем как вопросительный знак, или о Платоновской \textit{мании} поэзии, или о софистах, выступающих против истинного логоса лексическими уловками, или об Эдипе, прозревшем только став слепым, прикованном к своей зловещей судьбе, изолировавшей его и заклеймившей бездомным (άπολις).\footnote{Фигура \textit{аполиса} вырисовывается на фоне полиса, и Софокл сравнивает его с ύψίπολις, ''самым выдающимся в городе'' (\textit{Антигона} 5.370). Человек-\textit{аполис} соответствует ''лицу без гражданства'' в нынешнем понимании, но в нем также виден аналог \textit{атопоса}: ''лишенный места'', ''неуместный'', ''странный''. Alpha privativum в обоих случаях можно понимать как не только отрицание, но и выражение отказа, отстранения, выхода.} Такие маргинальные фигуры встречаются и в других местах: приверженность текстам в общепринятой части мистического учения обычно сопровождается ответвлениями, отбрасывающими тень неортодоксальности и анархии --- это касается и еврейской и исламской, и христианской традиций. Даже Святой Павел отмежевывается от риторики, в которой γλώσσα или \textit{lingua} расслабляют язык до такой степени, что позволяют ему ускользнуть от индивидуального или общественного контроля. Высказывание, растворяющееся в космосе сказанного --- заменяется ничего не говорящим и потому бессмысленным. Обаяние, противопоставляемое Максом Вебером повседневной жизни институтов, пронизывает все и показывает, что любая нормальность --- даже космологическая, теологическая или космополитически моделируемая --- упускает нечто, находящее затем выражение в некой безумной бахроме аномалий.

\section{Принадлежность, чуждость, контингентность}

Глобальный порядок немыслим без разрыва, если только место, из которого разворачивается целое, не превращается в конкретное целое. Модерн --- движение, ставящее под сомнение видение целого. Оно пресекает самообман космотеороса, считающего себя частью созерцаемой им сцены, и мнимый порядок оказывается результатом процесса установления порядка [\textit{Ordnungsstiftung}]. Здесь два ключевых открытия: (1) \textit{я}, говорящее ''я'' до того, как будет названо ''субъект'' и взрыв структуры отношений в целом из-за этой самореферентности и (2) радикальная \textit{контингентность}, которая не только занимает открытые пространства в порядке, но и влияет на сам порядок. Такой порядок может не только выродиться в беспорядок, но и стать другим порядком, быть отличным от самого себя. В cogito Декарта есть не только вершина мысли, но и возможность создания Богом другой математики. Порядки, в которых мы движемся, оказываются лишь потенциальными. И странный протагонист \textit{Человека без свойств} Музиля подозревает, что даже Бог предпочитает говорить о мире лишь потенциально: ''ибо Бог создает мир и в то же время думает, что он мог бы быть другим'' (Музиль, 1961, 15) и связывает это с тезисом ''действительно любящий свою страну не считает ее лучшей''. Решающий фактор здесь --- что мотивы субъектности и рациональности выступают как двойной мотив. Субъект модерна больше не олицетворяет лишь рациональность и ищет свое место, которого нет.

Если всерьез --- во всей его радикальности --- принять этот двойной мотив, то проблема проведения границ изменится: \textit{я}, живущее в своей собственной сфере и своему собственному культурному порядку соответствующее даже в языке и чувствах, уже не будет сводимо к чему-то всегда одному и отдельному от другого в контексте универсальности. Самость и самопринадлежност --- результат проведения границ, разделяющих внутреннее от внешнего и так принимающих формы \textit{включения} и \textit{исключения}. Чувство принадлежности возникает, когда что-то выходит из него, и это уходящее мы принимаем чуждым и неоднородным. Такое разделение своего и чуждого осуществляется без участия третьей стороны и составляет измерение, отдельное от ''то же / другое'' (которое опирается на диалектически созданное целое). Выражаясь языком западной философии, другой (έτερον) и чуждый (ξένον) --- не одно и то же. Чуждость гостя --- в том числе Незнакомца из Элеи, появляющегося в \textit{Софисте} Платона,--- иного языка или культуры не сводится к тому что нечто кажется другим. Дерево и бетон отличаются друг от друга, но никто не станет говорить, что они чужды друг другу. Чуждость значит, что у \textit{я} (\textit{ipse}) есть сфера принадлежности, собственного бытия, и что это \textit{я} не то же что \textit{то же} (\textit{idem}), которое различимо лишь третьей стороной.

\section{Компромиссы модерна}

Ницше пишет: ''Важно избавиться от Вселенной, от единства ... разрушить Вселенную, разучиться уважать ее'' (KSA 12:317). И едва ли модерн серьезно относится к последствиям разграничения порядков. До сих пор мы стремимся к компромиссам через взаимное смешение \textit{я} и \textit{того же} и \textit{другого} и \textit{чуждого}. Двойственность оценки потенциала нового порядка видна и в конфликте модерна и постмодерна. Взять хотя бы \textit{я}, которое говорит ''я'' и так отделяет себя от целого. ''Я'' буквально значит исключение. Прямое соединение картезианского эго с эгоцентричностью было бы чрезмерным упрощением, ведь пренебрегло бы беспокойством, проистекающим из этого радикального самоанализа и побуждающим философа искать незыблемый фундамент. Поиск обречен на провал, но что им движет? Даже если ''\textit{Кто} я?'' слишком быстро соскальзывает к ''\textit{Что} я?'', этот вопрос, замыкая вопрошающего и вопрошаемое, уже содержит след ''Je est un autre'' Артюра Рембо.

Как бы то ни было, открытие \textit{я}, чья самость отделяет его от структуры отношений в природном и социальном мире, ослабевает, когда его превращают в лишь \textit{частность}. В конце концов мы бы столкнулись только с конфликтом \textit{индивидуализма} и \textit{холизма}, \textit{партикуляризма} и \textit{универсализма} --- конфликта все еще сохраняющегося но никогда не приводившего к значимым результатам. Если целое нельзя понять без частей, в которых оно выражено, а общие правила малопригодны без учета обстоятельств частностей, к которым применяются, то конфликт этот приводит к великому объединению: аристотеликов --- с кантианцами, герменевтов --- с универсальными прагматиками. Но гармония нарушается, если вспомнить, что элементы ''Я-Ты-Здесь-Сейчас'' не принадлежат к общим концептуальным категориям --- скорее, эти индексальные или окказиональные выражения являются указателями на место произнесения речи, место развертывания поля опыта, языка и действия, прежде чем оно само займет какое-либо место. Место, о котором говорится, не совпадает с местом произнесения речи, и то же касается момента времени. Даже как некто, говорящий ''я'', я не являюсь счетным элементом класса или частью целого. ''Я'' произношения, будучи явно отлично от ''я'' содержания произнесенного --- не исчислимая вещь, о которой можно было бы говорить во множественном числе, и то же касается ''ты''. В древнегреческой мысли ''я'' и ''мы'' не играют тематической роли за исключением понятия ''для нас'' (πρός ήμας). Взрывная сила я-речи также теряется, если сводить ''я'' к общей я-функции. Здесь нам нужно защитить Декарта от его трансцендентально-филосовских наследников.

То же затруднение касается и второго --- приписывающего всем порядкам неустранимую контингентность --- открытия. Если рассматривать \textit{способность быть другим} или \textit{способность действовать по-другому} как большую или меньшую степень \textit{произвольности}, то антиподом полагается \textit{необходимость}. Касательно претензий на легитимность, это ведет к еще одному конфликту --- релятивизма и универсализма, оружие которых сейчас уже столь же изношено. Здесь стоит еще раз напомнить: установления порядков с их легитимностью --- включая генеалогию истинного и ложного или добра и зла --- не верно ни относительно, ни абсолютно. Оно никак не верно, потому что существование бинарных стандартов не подчинено самим этим стандартам, если только их генезис опять не скрыт и соответствующая оппозиция не гипостазирована. У каждого порядка есть свое слепое пятно --- что-то неупорядоченное, что не просто нехватка. Это относится как к моральным, так и к когнитивным и эстетическим порядкам. Поэтому современные авторы так упорно борются с \textit{инципитами} своих романов: с самого начала и писатель, и читатель попадают в ловушку готового порядка.\footnote{См. \textit{Феноменология внимания} (\textit{Phänomenologie der Aufmerksamkeit}), глава 3 \textit{Неисчислимое} (\textit{Unerzählbares}).} Этот порядок постанавляет: ''есть порядки'', где ''есть'' не поддается никакому обоснованию, являясь базовым допущением всякой попытки обоснования --- нулевым уровнем. То, что сейчас называется постмодернизмом, может значить, среди прочего, что некоторые поля сражений утратили свое значение --- но не гарантирует, что проблемы, проистекающие из модерновой субъектности и рациональности, не возникнут вновь уже в другом месте.

\section{Парадоксы самоограничения}

Парадоксы нынче в моде. Парадокс --- παρά δόξαν --- буквально мнение, расходящиеся с преобладающим. Как и в случае других пара- (паразитизм, парафизика, паралогизм, парапсихология), понятие парадокса предполагает нормальное положение вещей. Покуда парадоксы отклоняются лишь от преобладающих мнений, к ним нет возражений. Совсем другое дело, когда предположение вступает в противоречие с собственными предпосылками, усиливая одно предположение за счет другого, и наоборот. Для сохранения предположений значение парадоксов обычно преуменьшают, как в случае с парадоксом лживого критянина. Конечно, различение объектного и метаязыка спасает высказывание --- его событие теперь не часть его содержания. Но такое методологическое решение основано на вышеупомянутом ограничении, приходящим извне и требующем различения между двумя уровнями речи. Анализ и решение проблемы находятся на уровне говоримого, что значит пренебрежение возможностью высказывания и само-высказывания --- сказанного, но не явленного в речи. Самореферентность есть, например, в ''Я обещаю вам, что ...'', из-за нее критянин вовлекает себя во все, что говорит. Кроме того, хорошо известное различие между плоскостью содержания и плоскостью отношений из теории коммуникации Пола Вацлавика не может быть прямо применено \textit{к} речи, потому что всегда присутствует \textit{в} ней.

Нечто подобное происходит и при самоограничении, когда ограничивающее себя является результатом ограничения, а не некой третьей стороной, отделяющей одно от другого, не будучи в них вовлеченной. Эксперт по грибам ни ядовит, ни съедобен, и судья не имеет отношения к деяниям обвиняемого и ущербу истца. Но что происходит, когда некто извиняется? Легко вспомнить случаи, когда мы сталкивались с чем-то странным, от чего \textit{сами} ограждаются --- без третьей стороны: разница между близостью и удаленностью как разница между ''здесь'' и ''тогда'', ''сейчас'' и ''давным-давно'', чередующиеся миры бодрствования и сна, сверхъестественные ''беседы с умершими'', переносящие из одного мира в другой через провал в памяти в романе Хуана Рульфо ''Педро Парамо''\footnote{Интересный взгляд на мотив чуждого в современной литературе и теории литературы см. в \textit{Мексика за гранью одиночества --- попытка межкультурного анализа} (\textit{Mexiko jenseits der Einsamkeit -- Versuch einer interkulturellen Analyse}, 1994).}, отношения между мужчиной и женщиной, ребенком и взрослым, социальная изоляция из-за классовой принадлежности, профессии или культуры, наконец, границы нормы, отделяющие здоровых от больных, невиновных от преступников, ортодоксов от еретиков, своих от чужих. Во всех этих случаях относящий себя к одной из сторон оказывается вовне. Такие перемены, как засыпание и пробуждение, болезнь и выздоровление, старость, обращение в другую веру значат не новый вид от третьего лица, а становление чем-то другим. Пруст показал, в какой степени пробуждение зависит от чисто телесной самоориентации, его \textit{Исследования} начинаются на грани сна и бодрствования. Такие переживания --- когда свое сталкивается с чуждым --- преимущественно переживания пересечения границ и, как станет видно, сопровождаются удвоением движения. Двойная природа ограничения, когда, с одной стороны, некто отделен от другого, а с другой --- определенные области разграничены, допускает дальнейшую дифференциацию.

Первый аспект касается самореферентности действия по ограничению. Действующая граница не просто находится или проводится. Иллюстративно проведенная на доске линия здесь вводит в заблуждения. Акт проведения границы, когда некто отделяет себя от другого, нельзя ни увидить, ни потрогать --- он схватывается только через свой след. Его можно сравнить с заключением контракта --- действием, не становящимся частью самого контракта, но являемым косвенно, в изменении обязанностей. Процесс проведения границы происходит в нулевой точке, не находящейся ни здесь, ни там. Действующая граница потому --- ни определенная вещь, ни ничто, ведь без нее не было бы ни того, ни другого --- ни меня, ни других. \textit{Самореферентность проведения границ заключается в их самоустранении.}

Второй аспект касается самореферентности, проявляющейся в \textit{само} самоограничения и приводящей к внутреннему и внешнему ограничению. Эта самость --- ни завуалированный субъект собственного акта ограничения, ни объективный результат чуждого акта ограничения: она как бы выскакивает в самом акте ограничения, проявляясь как полость, \textit{как внутреннее, отделяющее себя от внешнего} и так создающее \textit{предпочтение в этой разнице}. Формально это значит: то, \textit{что} отделяет себя, отмечается, тогда как то, \textit{от чего} оно отделяется, остается неотмеченным. В этом неустранимая асимметрия, без которой не было бы себя, способного принять точку зрения другого или третьей стороны. И вышеупомянутые примеры иллюстрируют эту --- присущую контрасту между своим и чуждым --- односторонность.

\end{document}
